\section{Conclusions and Future Research Directions}
\label{sec:conclusions}

This paper has presented S4D-TAM (Semantic 4D Token Attention Map), a novel autonomous navigation architecture for unmanned aerial vehicles operating in GNSS-degraded environments. The foundational contribution is the persistent 4D semantic token map, unifying localization, semantic scene understanding, spatiotemporal occupancy forecasting, risk estimation, and motion planning into a single coherent spatial representation.

The key conceptual contributions include:
\begin{enumerate}
  \item \textbf{4D Semantic Token Model}: Compactly integrating spatial geometry, continuous semantics, dynamic state, observation history, and calibrated uncertainty.
  \item \textbf{Hierarchical Multi-Layer Architecture}: Unifying multi-sensor fusion, SLAM front-end, tokenization, transformer attention, memory persistence, reference map alignment, probabilistic forecasting, risk-aware trajectory planning, and active perception.
  \item \textbf{Spatiotemporal Attention Mechanism}: Gating attention by spatial proximity, temporal decay, and uncertainty bounds to guarantee robust state estimation.
  \item \textbf{Proactive Occlusion and Risk Forecasting}: Inferring hazard probabilities in unobserved spatial volumes conditioned on semantic topology and dynamic flows.
  \item \textbf{Multi-Objective Risk-Aware Planning}: Optimizing flight trajectories against safety risk, path length, energy expenditure, and spatial uncertainty.
  \item \textbf{Information-Theoretic Active Perception}: Dynamically allocating sensor modalities and processing rates based on mutual information gain and energy constraints.
\end{enumerate}

Seven preregistered mechanism hypotheses (H1--H7) have been defined to causally isolate the impact of individual architectural features from complete-system baseline comparisons against ORB-SLAM3, VINS-Mono, FAST-LIO2, and LIO-SAM.

Future research directions will concentrate on:
\begin{enumerate}
  \item Executing the prospective two-level validation study across TartanAir, Blackbird, MARSIM, and AeroVerse datasets, progressing from offline benchmarks to SIL, HIL, and field flight tests.
  \item Embedded model optimization through post-training quantization, structured token pruning, and hardware-accelerated attention backends.
  \item Expanding multi-sensor modalities to millimeter-wave radar and ultrasound, and evaluating multi-UAV cooperative mapping teams.
  \item Integrating multimodal large language models (LLMs) for high-level semantic reasoning and natural language mission command.
\end{enumerate}
